% ===========================================================================
%  Annexe — Mathématiques financières, II. Annuités et emprunts indivis
%  Hors des quatre composantes chiffrées du PE 2026 ; RAPE, période 4 ;
%  exercice 3 de chaque sujet de baccalauréat OSE examiné.
% ===========================================================================
\chapter{Annuités et emprunts indivis}
\markboth{Mathématiques financières}{Annuités et emprunts indivis}

Second volet de cette annexe : calculer la valeur acquise et la valeur
actuelle d'une suite d'annuités constantes, établir un tableau
d'amortissement d'un emprunt indivis, et comparer les modes de remboursement
par annuités constantes et par amortissements constants.

Une coopérative qui épargne chaque année la même somme, ou un emprunteur qui
rembourse chaque année la même mensualité : dans les deux cas, on retrouve
une \textbf{suite d'annuités constantes} — un versement identique, répété à
intervalles réguliers. Ce second volet referme l'annexe consacrée aux
mathématiques financières en répondant à la question la plus concrète :
combien épargne-t-on au bout de $n$ ans de versements réguliers, et combien
emprunte-t-on quand on promet de rembourser une somme fixe chaque année ?

Les deux questions sont les deux faces d'une même formule : la première
regarde vers l'avenir — la \textbf{valeur acquise} — la seconde regarde vers
le présent — la \textbf{valeur actuelle}. C'est cette seconde lecture qui
fonde le calcul d'un emprunt et de son tableau d'amortissement, avec lequel
se termine cette annexe.

% ---------------------------------------------------------------------------
\begin{revision}
Premier volet de cette annexe et suites géométriques.

\begin{enumerate}[label=\textbf{\arabic*.}, leftmargin=*, itemsep=0.35em]
  \item Rappeler la somme des $n$ premiers termes d'une suite géométrique de
        premier terme $u_{0}$ et de raison $q \neq 1$.
  \item Calculer $(1{,}08)^{3}$.
  \item Rappeler la formule de la valeur acquise d'un capital $C$ à intérêt
        composé au taux $t$ pendant $n$ années.
  \item Que vaut $(1+t)^{-n}$ lorsque $n$ augmente indéfiniment, pour
        $t>0$ ?
  \item Résoudre $\dfrac{1-x}{0{,}09} = 3$ d'inconnue $x$.
\end{enumerate}
\end{revision}

\begin{automatismes}
Sans calculatrice, sans brouillon.

\begin{enumerate}[label=\textbf{\alph*)}, leftmargin=*, itemsep=0.25em]
  \item $(1+t)^{0}$
  \item Somme de $5$ annuités de $10\,000$~Ar sans capitalisation
  \item $(1{,}1)^{2}$
  \item $(1+t)^{-1}$ pour $t = 0{,}1$
  \item Capital restant dû après remboursement total d'un emprunt
  \item Un emprunt de $200\,000$~Ar remboursé en $4$ amortissements
        constants : montant de chaque amortissement ?
  \item Intérêt de la première année d'un emprunt de $200\,000$~Ar au taux
        de $10\,\%$
  \item Annuité $=$ amortissement $+$ \dots
\end{enumerate}
\end{automatismes}

\begin{corrige}[automatismes]
a) $1$ \quad b) $50\,000$ \quad c) $1{,}21$ \quad d) $\approx 0{,}909$
\quad e) $0$ \quad f) $50\,000$ \quad g) $20\,000$ \quad h) intérêt.
\end{corrige}

% ===========================================================================
\section{Valeur acquise d'une suite d'annuités constantes}

\begin{definition}[annuité constante]
Une \textbf{annuité} est un versement effectué à intervalles de temps
réguliers, en général chaque année, en fin de période. Elle est
\textbf{constante} lorsque tous les versements ont le même montant $a$.
\end{definition}

\begin{theoreme}[valeur acquise d'une suite d'annuités constantes]
La \textbf{valeur acquise} à la date du dernier versement, de $n$ annuités
constantes $a$, versées en fin de période au taux $t$, vaut
\[
  V_{n} = a\,\frac{(1+t)^{n}-1}{t} .
\]
\end{theoreme}

\begin{demonstration}
Le premier versement, effectué à la fin de la première période, produit des
intérêts composés pendant les $n-1$ périodes restantes et vaut donc
$a(1+t)^{n-1}$ à la date finale ; le second, $a(1+t)^{n-2}$ ; et ainsi de
suite jusqu'au dernier, versé sans capitalisation. La valeur acquise totale
est la somme
\[
  V_{n} = a+a(1+t)+\cdots+a(1+t)^{n-1}
  = a\,\frac{(1+t)^{n}-1}{(1+t)-1} = a\,\frac{(1+t)^{n}-1}{t},
\]
somme des $n$ premiers termes d'une suite géométrique de raison $1+t$.
\end{demonstration}

\begin{exemple}[valeur acquise]
Une coopérative verse, en fin d'année, cinq annuités constantes de
$80\,000$~Ar sur un compte rémunéré à $7\,\%$.
\[
  V_{5} = 80\,000 \times \frac{(1{,}07)^{5}-1}{0{,}07}
  \approx 80\,000 \times 5{,}750\,74 \approx 460\,059{,}1 \text{ Ar}.
\]
\end{exemple}

% ===========================================================================
\section{Valeur actuelle d'une suite d'annuités constantes}

\begin{theoreme}[valeur actuelle d'une suite d'annuités constantes]
La \textbf{valeur actuelle}, une période avant le premier versement, de
$n$ annuités constantes $a$, au taux $t$, vaut
\[
  V_{0} = a\,\frac{1-(1+t)^{-n}}{t} .
\]
\end{theoreme}

\begin{remarque}
On passe de $V_{n}$ à $V_{0}$ en actualisant : $V_{0} = V_{n}(1+t)^{-n}$,
ce que l'on retrouve directement en factorisant $(1+t)^{n}$ dans la formule
précédente. Cette valeur actuelle est aussi le montant qu'un organisme
prêteur accepte de verser aujourd'hui, en échange de $n$ remboursements
futurs de $a$ : c'est le principe même de l'\textbf{emprunt indivis}, étudié
à la section suivante.
\end{remarque}

\begin{exemple}[valeur actuelle — montant d'un emprunt]
Une coopérative rembourse un emprunt par $4$ annuités constantes de
$150\,000$~Ar, au taux annuel de $9\,\%$.
\[
  V_{0} = 150\,000 \times \frac{1-(1{,}09)^{-4}}{0{,}09}
  \approx 150\,000 \times 3{,}239\,73 \approx 485\,959{,}1 \text{ Ar}.
\]
Le montant emprunté était donc d'environ $485\,959$~Ar.
\end{exemple}

\begin{entrainement}[valeur acquise et valeur actuelle]
\exo[\niveau{1}] Calculer la valeur acquise de $6$ annuités constantes de
$50\,000$~Ar, au taux de $6\,\%$.

\exo[\niveau{2}] Déterminer le montant de chaque annuité constante permettant
d'obtenir une valeur acquise de $400\,000$~Ar en $5$ ans, au taux de $8\,\%$.

\exo[\niveau{2}] Calculer la valeur actuelle de $5$ annuités constantes de
$90\,000$~Ar, au taux de $7\,\%$.

\exo[\niveau{3}] Un épargnant verse $10$ annuités constantes de $60\,000$~Ar
au taux de $5\,\%$, puis laisse la somme obtenue fructifier encore $3$ ans au
même taux, sans nouveau versement. Calculer le capital final.
\end{entrainement}

% ===========================================================================
\section{Emprunt indivis et tableau d'amortissement}

\begin{definition}[emprunt indivis, amortissement]
Un \textbf{emprunt indivis} est un emprunt contracté auprès d'un seul
prêteur, remboursé par une suite de versements annuels appelés
\textbf{annuités}. Chaque annuité se décompose en un \textbf{intérêt}, calculé
sur le capital restant dû, et un \textbf{amortissement}, qui rembourse une
partie de ce capital :
\[
  \text{annuité} = \text{intérêt}+\text{amortissement} .
\]
\end{definition}

\begin{propriete}[relations du tableau d'amortissement]
Soit $K_{0}$ le capital emprunté, $t$ le taux annuel, $a$ l'annuité (ici
supposée constante) et $K_{k}$ le capital restant dû après le $k$-ième
versement. Pour tout $k$,
\[
  i_{k} = t\,K_{k-1}, \qquad
  m_{k} = a-i_{k}, \qquad
  K_{k} = K_{k-1}-m_{k},
\]
où $i_{k}$ et $m_{k}$ désignent respectivement l'intérêt et l'amortissement
de la $k$-ième période, et $K_{n} = 0$ après le dernier versement.
\end{propriete}

\begin{theoreme}[annuité constante d'un emprunt]
Pour rembourser un capital $K_{0}$ en $n$ annuités constantes au taux $t$,
chaque annuité vaut
\[
  a = \frac{K_{0}\,t}{1-(1+t)^{-n}} .
\]
\end{theoreme}

\begin{demonstration}
Le capital emprunté est la valeur actuelle des $n$ annuités qui le
remboursent : $K_{0} = a\,\dfrac{1-(1+t)^{-n}}{t}$, d'où l'expression de
$a$ en isolant cette inconnue.
\end{demonstration}

\begin{methode}{dresser un tableau d'amortissement}
\begin{enumerate}[label=\textbf{\arabic*.}, leftmargin=*, itemsep=0.15em]
  \item Calculer l'annuité constante $a$ à l'aide de la formule précédente.
  \item Pour chaque période, calculer l'intérêt $i_{k} = t\,K_{k-1}$ sur le
        capital restant en début de période.
  \item En déduire l'amortissement $m_{k} = a-i_{k}$, puis le nouveau
        capital restant $K_{k} = K_{k-1}-m_{k}$.
  \item Vérifier qu'après le dernier versement, $K_{n}$ est nul — à un
        arrondi près, dû à l'arrondi de $a$.
\end{enumerate}
\end{methode}

\begin{exemple}[un tableau d'amortissement complet]
Un emprunt de $300\,000$~Ar est remboursé en $3$ annuités constantes, au
taux de $8\,\%$. L'annuité vaut
$a = \dfrac{300\,000 \times 0{,}08}{1-(1{,}08)^{-3}} \approx 116\,410$~Ar.
\[
  \begin{array}{c|cccc}
    \text{Période} & \text{Capital dû (début)} & \text{Intérêt} &
    \text{Amortissement} & \text{Capital dû (fin)} \\
    \hline
    1 & 300\,000 & 24\,000 & 92\,410 & 207\,590 \\
    2 & 207\,590 & 16\,607{,}2 & 99\,802{,}8 & 107\,787{,}2 \\
    3 & 107\,787{,}2 & 8\,622{,}98 & 107\,787{,}02 & \approx 0
  \end{array}
\]
Le capital restant s'annule bien en fin de tableau, à quelques centimes près
dus à l'arrondi de l'annuité — ce léger écart, normal, doit être signalé
plutôt que caché.
\end{exemple}

\begin{remarque}[amortissements constants]
Un autre mode de remboursement, plus simple, consiste à fixer non
l'annuité mais l'\textbf{amortissement}, à raison de $m = \dfrac{K_{0}}{n}$
identique chaque période. L'intérêt, calculé sur un capital restant qui
diminue chaque fois de la même somme, décroît alors régulièrement, et
l'annuité — intérêt plus amortissement — décroît elle aussi d'une période à
l'autre, contrairement au cas précédent où elle restait fixe.
\end{remarque}

\begin{entrainement}[emprunt et amortissement]
\exo[\niveau{1}] Un emprunt de $250\,000$~Ar est remboursé par $4$
amortissements constants. Calculer le montant de chaque amortissement.

\exo[\niveau{2}] Un emprunt de $600\,000$~Ar, au taux de $9\,\%$, est
remboursé en $5$ annuités constantes. Calculer le montant de l'annuité.

\exo[\niveau{2}] Avec les données précédentes, calculer l'intérêt et
l'amortissement de la première période, puis le capital restant dû après ce
premier remboursement.

\exo[\niveau{3}] Montrer que, dans un tableau d'amortissement à annuités
constantes, la suite des amortissements $\left(m_{k}\right)$ est
géométrique de raison $1+t$.
\end{entrainement}

% ===========================================================================
\begin{annales}
Les annuités ferment l'exercice de mathématiques financières, en général sur
un scénario de placement ou d'emprunt à taux constant, parfois complété par
la recherche du taux lui-même. L'énoncé qui suit est repris d'un sujet
réellement posé.

\exoann{Baccalauréat de l'Enseignement Général, Madagascar, série OSE, option OSE, session 2026 — exercice 3, partie III}
On verse $6$ annuités constantes au taux annuel de $10\,\%$, et l'on obtient
une valeur acquise de $771\,561{,}00$~Ar.
\begin{enumerate}[label=\textbf{\alph*)}, leftmargin=*, itemsep=0.15em]
  \item Déterminer le montant de chaque annuité.
\end{enumerate}

On verse à présent $6$ annuités constantes de $100\,000$~Ar chacune, qui
produisent une valeur acquise de $811\,518{,}90$~Ar.
\begin{enumerate}[label=\textbf{\alph*)}, leftmargin=*, itemsep=0.15em]
  \setcounter{enumii}{1}
  \item Déterminer le taux annuel $t$ de ce second placement.
\end{enumerate}

\exoann[emprunt remboursé par une suite]{Baccalauréat, France, série ES, Pondichéry, 21 avril 2016 — exercice 4 (extrait, sans l'algorithme)}
En janvier 2016, une personne achète un scooter coûtant $5\,700$ euros sans
apport, financé par un crédit à la consommation de $5\,700$ euros au taux
mensuel de $1{,}5\,\%$. La mensualité, fixée à $300$ euros, est versée le
$25$ de chaque mois ; ainsi le capital restant dû augmente de $1{,}5\,\%$
puis diminue de $300$ euros. Le premier versement a lieu le 25 février 2016.
On note $u_{n}$ le capital restant dû, en euros, juste après la $n$-ième
mensualité, avec $u_{0} = 5\,700$.
\begin{enumerate}[label=\textbf{\arabic*.}, leftmargin=*, itemsep=0.2em]
  \item Montrer que $u_{1} = 5\,485{,}50$ euros, puis calculer $u_{2}$.
  \item On admet que, pour tout entier naturel $n$,
        $u_{n+1} = 1{,}015\,u_{n}-300$. On pose $v_{n} = u_{n}-20\,000$.
        \begin{enumerate}[label=\textbf{\alph*)}, leftmargin=1.2em]
          \item Montrer que $(v_{n})$ est géométrique de raison $1{,}015$, et
                en déduire que, pour tout $n$,
                $u_{n} = 20\,000-14\,300\times1{,}015^{n}$.
          \item Vérifier qu'une valeur approchée du capital restant dû au
                26 avril 2017 (juste après la $15^{\text{e}}$ mensualité) est
                $2\,121{,}68$ euros.
        \end{enumerate}
  \item Déterminer le nombre de mensualités nécessaires pour rembourser
        intégralement le prêt, puis le montant de la dernière mensualité
        (inférieur à $300$ euros, puisqu'elle solde le capital restant).
        Quel est alors le coût total de l'achat du scooter ?
\end{enumerate}
\end{annales}

\begin{corrige}[annales]
\textbf{a)} La valeur acquise de $6$ annuités constantes $a$ au taux de
$10\,\%$ vaut
$V_{6} = a\,\dfrac{(1{,}1)^{6}-1}{0{,}1}$. Or $(1{,}1)^{6} = 1{,}771\,561$,
donc
\[
  V_{6} = a \times \frac{0{,}771\,561}{0{,}1} = a \times 7{,}715\,61 .
\]
En résolvant $771\,561 = 7{,}715\,61\,a$, on trouve
$a = \dfrac{771\,561}{7{,}715\,61} = 100\,000$~Ar.

\textbf{b)} Avec $a = 100\,000$ et $V_{6} = 811\,518{,}90$, l'équation
$100\,000 \times \dfrac{(1+t)^{6}-1}{t} = 811\,518{,}90$ doit être résolue
en $t$. En testant $t = 0{,}12$ :
$(1{,}12)^{6} = 1{,}973\,822\,685$, d'où
\[
  100\,000 \times \frac{1{,}973\,822\,685-1}{0{,}12}
  = 100\,000 \times \frac{0{,}973\,822\,685}{0{,}12}
  = 100\,000 \times 8{,}115\,189
  = 811\,518{,}90 .
\]
La valeur trouvée coïncide exactement avec celle de l'énoncé : le taux
annuel est donc $t = 12\,\%$.

\textbf{Exercice suivant — 1.} Chaque mois, le capital restant dû augmente
de $1{,}5\,\%$ puis diminue de la mensualité de $300$ euros :
$u_{1} = 1{,}015\times5\,700-300 = 5\,785{,}50-300 = 5\,485{,}50$ euros, puis
$u_{2} = 1{,}015\times5\,485{,}50-300 = 5\,567{,}78-300 = 5\,267{,}78$ euros.

\textbf{2. a)} $v_{n+1} = u_{n+1}-20\,000 = 1{,}015\,u_{n}-20\,300
= 1{,}015\,(u_{n}-20\,000) = 1{,}015\,v_{n}$ : $(v_{n})$ est géométrique de
raison $1{,}015$ et de premier terme $v_{0} = 5\,700-20\,000 = -14\,300$.
Donc $v_{n} = -14\,300\times1{,}015^{n}$ et
$u_{n} = 20\,000-14\,300\times1{,}015^{n}$.
\textbf{b)} Le 26 avril 2017 tombe juste après la $15^{\text{e}}$
mensualité (versée de février 2016 à avril 2017 inclus). Comme
$1{,}015^{15} \approx 1{,}250232$,
\[
  u_{15} \approx 20\,000-14\,300\times1{,}250232 \approx
  20\,000-17\,878{,}32 \approx 2\,121{,}68 \text{ euros.}
\]

\textbf{3.} Le prêt est soldé au premier rang $n$ tel que $u_{n}\leq0$, soit
$1{,}015^{n} \geq \dfrac{20\,000}{14\,300} \approx 1{,}3986$, c'est-à-dire
$n \geq \dfrac{\ln(1{,}3986)}{\ln(1{,}015)} \approx 22{,}5$. Le remboursement
demande donc $23$ mensualités : les $22$ premières de $300$ euros, la
dernière soldant le reste. Après la $22^{\text{e}}$ mensualité,
$u_{22} \approx 20\,000-14\,300\times1{,}015^{22} \approx 163{,}94$ euros
restent dus ; avec l'intérêt du dernier mois
($0{,}015\times163{,}94\approx2{,}46$), la dernière mensualité vaut environ
$166{,}40$ euros. Le coût total de l'achat s'élève donc à
$22\times300+166{,}40 = 6\,766{,}40$ euros, soit un coût du crédit d'environ
$1\,066{,}40$ euros au-dessus du prix affiché du scooter.
\end{corrige}

% ===========================================================================
\begin{jecherche}[l'emprunt du grenier communautaire]
Une association villageoise emprunte $500\,000$~Ar pour construire un
grenier communautaire destiné au stockage du riz. Elle rembourse cet emprunt
par $4$ annuités constantes, au taux annuel de $9\,\%$.

\begin{enumerate}[label=\textbf{\arabic*.}, leftmargin=*, itemsep=0.3em]
  \item Calculer le montant de l'annuité constante $a$.
  \item Dresser les deux premières lignes du tableau d'amortissement
        (période, capital dû en début de période, intérêt, amortissement,
        capital dû en fin de période).
  \item Sans calculer les périodes suivantes, expliquer pourquoi la somme
        des quatre amortissements doit nécessairement valoir $500\,000$~Ar.
  \item Calculer le coût total de l'emprunt — c'est-à-dire la somme des
        quatre annuités moins le capital emprunté.
  \item L'association envisage un autre mode de remboursement, par
        \textbf{amortissements constants} de $125\,000$~Ar chacun. Calculer,
        pour ce second mode, le total des intérêts versés sur les quatre
        années.
  \item Comparer les deux modes de remboursement du point de vue du coût
        total, et expliquer d'où vient la différence.
\end{enumerate}
\end{jecherche}

\begin{corrige}[l'emprunt du grenier communautaire]
\textbf{1.} $a = \dfrac{500\,000 \times 0{,}09}{1-(1{,}09)^{-4}}
\approx \dfrac{45\,000}{0{,}291\,58} \approx 154\,334$~Ar.

\textbf{2.}
\[
  \begin{array}{c|cccc}
    \text{Période} & \text{Capital dû (début)} & \text{Intérêt} &
    \text{Amortissement} & \text{Capital dû (fin)} \\
    \hline
    1 & 500\,000 & 45\,000 & 109\,334 & 390\,666 \\
    2 & 390\,666 & 35\,159{,}9 & 119\,174{,}1 & 271\,491{,}9
  \end{array}
\]

\textbf{3.} Chaque amortissement rembourse une partie du capital emprunté,
et le capital restant dû doit s'annuler après le dernier versement — c'est
la définition même d'un emprunt totalement remboursé. La somme de tous les
amortissements est donc nécessairement égale au capital initial
$K_{0} = 500\,000$~Ar, quel que soit le détail de leur répartition dans le
temps.

\textbf{4.} Le total des annuités vaut $4 \times 154\,334 = 617\,336$~Ar ;
le coût de l'emprunt est donc $617\,336-500\,000 = 117\,336$~Ar (ce montant
correspond, aux arrondis près, à la somme des quatre intérêts annuels).

\textbf{5.} Avec des amortissements constants de $125\,000$~Ar, le capital
restant dû vaut successivement $500\,000$, $375\,000$, $250\,000$ et
$125\,000$~Ar en début de chaque période, d'où les intérêts
\[
  45\,000, \quad 33\,750, \quad 22\,500, \quad 11\,250 \text{ Ar},
\]
dont la somme vaut $112\,500$~Ar.

\textbf{6.} Le remboursement par amortissements constants coûte moins cher
en intérêts totaux ($112\,500$~Ar contre environ $117\,336$~Ar) : en
imposant un amortissement identique dès la première période, le capital
restant dû diminue plus vite en moyenne que dans le cas des annuités
constantes, où les premiers amortissements sont plus faibles. Un capital
restant dû plus faible, en moyenne, produit moins d'intérêts cumulés — même
si, en contrepartie, la première annuité à amortissements constants est plus
lourde à supporter pour l'emprunteur.
\end{corrige}

% ===========================================================================
\begin{jemeteste}
Répondre par vrai ou faux, en justifiant en une ligne.

\begin{enumerate}[label=\textbf{\arabic*.}, leftmargin=*, itemsep=0.28em]
  \item Une annuité est toujours égale à un amortissement.
  \item Le capital restant dû s'annule après le dernier remboursement d'un
        emprunt entièrement amorti.
  \item Dans le remboursement par annuités constantes, l'intérêt annuel
        augmente au fil du temps.
  \item La valeur actuelle d'une suite d'annuités constantes correspond au
        montant qu'un emprunt de mêmes caractéristiques permettrait
        d'obtenir.
  \item Dans le remboursement par amortissements constants, l'annuité
        diminue au fil du temps.
  \item La somme des amortissements d'un emprunt totalement remboursé est
        toujours égale au capital emprunté.
\end{enumerate}
\end{jemeteste}

\begin{corrige}[je me teste]
\textbf{1.} Faux — l'annuité est la somme de l'intérêt et de
l'amortissement.
\textbf{2.} Vrai, par définition d'un emprunt totalement remboursé.
\textbf{3.} Faux — c'est l'inverse : le capital restant dû diminue, donc
l'intérêt diminue lui aussi.
\textbf{4.} Vrai — c'est exactement le principe de l'emprunt indivis.
\textbf{5.} Vrai — l'intérêt décroît d'une période à l'autre tandis que
l'amortissement reste fixe.
\textbf{6.} Vrai.
\end{corrige}

% ===========================================================================
\begin{commeaubac}
\bmtsource{Exercice original, au format de l'épreuve — série OSE}
Une coopérative agricole emprunte pour l'achat d'une décortiqueuse à riz.
L'emprunt est remboursé par $5$ annuités constantes au taux annuel de
$8\,\%$, et la valeur acquise de ces cinq annuités, à la date du dernier
versement, vaut $703\,992{,}12$~Ar.

\begin{enumerate}[label=\textbf{\arabic*.}, leftmargin=*, itemsep=0.25em]
  \item Calculer le montant de chaque annuité constante $a$.
        \hfill\textup{(1 pt)}
  \item En déduire le montant emprunté, c'est-à-dire la valeur actuelle de
        ces cinq annuités. \hfill\textup{(0,75 pt)}
  \item Calculer l'intérêt et l'amortissement de la première période, ainsi
        que le capital restant dû après ce premier remboursement.
        \hfill\textup{(0,75 pt)}
\end{enumerate}
\end{commeaubac}

\begin{corrige}[comme au baccalauréat]
\textbf{1.} $(1{,}08)^{5} = 1{,}469\,328$, donc
\[
  V_{5} = a \times \frac{1{,}469\,328-1}{0{,}08}
  = a \times 5{,}866\,60 = 703\,992{,}12,
\]
d'où $a = \dfrac{703\,992{,}12}{5{,}866\,60} = 120\,000$~Ar.

\textbf{2.} Le montant emprunté est la valeur actuelle des cinq annuités :
\[
  K_{0} = a \times \frac{1-(1{,}08)^{-5}}{0{,}08}
  \approx 120\,000 \times 3{,}992\,71 \approx 479\,125{,}50 \text{ Ar}.
\]

\textbf{3.} Le premier intérêt vaut $i_{1} = 0{,}08 \times 479\,125{,}50
\approx 38\,330{,}04$~Ar ; l'amortissement de la première période vaut donc
$m_{1} = 120\,000-38\,330{,}04 = 81\,669{,}96$~Ar, et le capital restant dû
après ce remboursement est
$K_{1} = 479\,125{,}50-81\,669{,}96 = 397\,455{,}54$~Ar.
\end{corrige}

% ===========================================================================
\begin{notepourlaclasse}
Ce second volet récompense les élèves qui maîtrisent déjà les suites
géométriques : la formule de la valeur acquise n'est rien d'autre qu'une
somme de termes géométriques, et le rappeler explicitement, plutôt que de la
présenter comme une formule isolée à mémoriser, la rend beaucoup plus
robuste en mémoire.

Le tableau d'amortissement doit être construit ligne après ligne, jamais
recopié d'une formule fermée pour chaque $K_{k}$ : c'est la répétition du
même petit calcul — intérêt, puis amortissement, puis nouveau capital — qui
ancre la logique du dispositif, bien plus qu'une formule générale peu
lisible.

Enfin, le sujet 2026 illustre une compétence rarement travaillée mais
explicitement attendue : retrouver un taux d'intérêt par tâtonnement
raisonné à partir d'une valeur acquise donnée, faute de pouvoir résoudre
l'équation directement. Faites pratiquer cette démarche — poser une valeur
plausible, vérifier, ajuster — comme une méthode à part entière, et non
comme un pis-aller.
\end{notepourlaclasse}
